\apendice{Especificación de diseño}

\section{Introducción}

Una vez definidos los objetivos y requisitos del sistema, en este apartado, se describe el diseño de la aplicación desarrollada. La especificación de diseño tiene como finalidad definir la estructura, organización e interacción de los distintos componentes que forman el sistema, estableciendo la base técnica necesaria para su implementación.

El objetivo principal de esta fase es trasladar los requisitos funcionales y no funcionales previamente definidos a una solución técnica coherente, modular y escalable, permitiendo garantizar la correcta gestión de la información, la estabilidad del sistema y el adecuado funcionamiento de todos los procesos implicados en la aplicación.

A lo largo de este apéndice se detallan los principales elementos que constituyen la arquitectura de PythIA. Para ello, se describe el diseño de los datos utilizados por la aplicación (ver sección~\hyperref[dis:dtos]{C.2}), la arquitectura general del sistema (ver sección~\hyperref[dis:arc]{C.3}), la organización de sus componentes software y la relación existente entre ellos.

Además, se incluyen aspectos relacionados con el diseño procedimental (ver sección~\hyperref[dis:proc]{C.4}) y el diseño de interfaces (ver sección~\hyperref[dis:int]{C.5}), describiendo las interfaces de usuario y la navegación del sistema.

Este diseño busca facilitar tanto el desarrollo y mantenimiento de la aplicación como su futura evolución, permitiendo incorporar nuevas funcionalidades, modelos y componentes con el menor impacto posible sobre la arquitectura existente.

\section{Diseño de datos} \label{dis:dtos}

Para llevar a cabo la aplicación se han definido las siguientes entidades principales (ver figura~\ref{fig:Diagramas/Entidades})\footnote{Esta figura se ha generado utilizando la aplicación DBeaver (\href{https://dbeaver.io/}{https://dbeaver.io/})}
\begin{itemize}
	\item \textbf{Usuario}: esta entidad representa a cada persona registrada en la aplicación. Distingue entre usuario normal y administrador, al cual se le permite realizar la gestión de documentos y usuarios. Esta entidad maneja un identificador (id) que es la clave primaria, un nombre, un \eng{email} que debe ser único, una contraseña que se almacena cifrada, el último inicio de sesión, el rol, el número de inicios de sesión realizados, el país, el idioma preferido, el modelo de lenguaje seleccionado por defecto y el modo visual de la interfaz.
	\item \textbf{Consulta}: esta entidad guarda las preguntas que realizan los usuarios, así como las respuestas que da el \eng{LLM} alimentado por el \textit{RAG}. Es la base para el historial de consultas. Cada consulta registra la pregunta formulada, la respuesta obtenida, los fragmentos documentales utilizados como contexto, el tiempo de respuesta, el dispositivo empleado para la ejecución y la fecha de realización de la consulta.
	\item \textbf{Documento}: esta entidad representa un PDF cargado en la aplicación. Su función es mantener la trazabilidad completa del ciclo de vida documental, desde su carga inicial hasta su conversión e indexación. Además de almacenar el nombre del fichero, la ruta y el tamaño, registra el \eng{hash}, el estado de procesamiento, el contenido Markdown generado y diversos metadatos asociados como el número de expediente o el tipo documental. En caso de que sucedan, también almacena los errores ocurridos durante el procesamiento. 
	\item \textbf{Chunk}: esta entidad representa cada trozo de texto en el que se dividen los documentos para poder generar los \textit{embeddings}. Cada fragmento almacena información sobre el documento al que pertenece, su posición dentro del mismo, el identificador asociado en la base de datos vectorial, el \eng{hash} de la versión del documento, el número de expediente, el tipo documental y otros metadatos necesarios para la recuperación semántica de información.
	\item \textbf{Embedding}: esta entidad almacena los metadatos asociados a la respresentación vectorial de un fragmento (\eng{chunks}). A diferencia del vector numérico propiamente dicho, que se almacena en Qdrant, esta entidad conserva información de trazabilidad necesaria para gestionar el proceso de indexación. Concretamente almacena el modelo utilizado para generar el \eng{embedding}, la dimensionalidad del vector, la métrica de similitud empleada y la fecha de creación.
	\item \textbf{ConsultaChunk}: esta entidad modela la relación entre las consultas y los \textit{chunks} recuperados durante la búsqueda semántica. Permite registrar qué fragmentos fueron utilizados para generar cada respuesta. Además de los identificadores de la consulta y del fragmento, almacena la puntuación de similitud obtenida durante la recuperación y la posición ocupada por cada fragmento dentro del \eng{ranking} de relevancia.
\end{itemize}

\imagendiagram{Diagramas/Entidades}{Diagrama Entidad/Relación de las entidades principales generado por DBeaver.}{1.1}

Además de las entidades principales del sistema, se han definido varias entidades destinadas a gestionar y monitorizar la ejecución de tareas realizadas de forma asíncrona (ver figura~\ref{fig:Diagramas/Entidades_asincronas})\footnote{Esta figura se ha generado utilizando la aplicación DBeaver (\href{https://dbeaver.io/}{https://dbeaver.io/})}. La mayoría de ellas heredan de la clase abstracta \textbf{State}, mientras que la entidad \textbf{RAGQueryState} implementa una estructura equivalente mediante la clase auxiliar \textbf{JobStateMixin}. En ambos casos se reutiliza un conjunto común de atributos y métodos destinados a registrar el estado de ejecución de una tarea:

\begin{itemize}
	\item \textbf{RAGQueryState}: esta entidad permite gestionar las consultas en ejecución de forma asíncrona. Almacena la pregunta enviada, el modelo seleccionado,  el estado de ejecución, los resultados obtenidos y los errores producidos (si los hubiese). 
	\item \textbf{MarkdownConversionState}: esta entidad registra el estado de los procesos de conversión de documentos PDF a Markdown. Permite conocer el proceso de la tarea, los mensajes informativos generados y los posibles errores.
	\item \textbf{VectorUpdateState}: esta entidad controla los procesos de actualización de la base de datos vectorial. Almacena información sobre el documento que se está procesando en cada momento y los posibles errores producidos.
	\item \textbf{WebScrapingState}: esta entidad registra el estado de los procesos de \eng{web scraping}. Esta entidad permite monitorizar la ejecución, conocer le progreso alcanzado y registrar incidencias durante la descarga de documentos.
	\item \textbf{RAGEvaluationState}: esta entidad almacena la información asociada a las evaluaciones automáticas del sistema \eng{RAG}. Registra el progreso de la evaluación, los resultados obtenidos y las rutas de los distintos archivos generados durante el proceso, incluyendo configuraciones, conjuntos de datos y métricas calculadas.
\end{itemize}

\imagendiagram{Diagramas/Entidades_asincronas}{Diagrama Entidad/Relación de las entidades asíncronas generado por DBeaver.}{1.1}


Las relaciones entre las entidades del sistema son (ver figura~\ref{fig:entidad_relacion}):

\begin{itemize}
	\item Un \textbf{usuario} puede realizar cero o muchas \textbf{consultas}, mientras que cada consulta pertenece a un único usuario. Esta relación permite mantener el historial de consultas asociado a cada cuenta.
	\item Un \textbf{usuario} puede tener cero o muchas instancias de \textbf{RAGQueryState}, utilizadas para gestionar consultas ejecutadas de forma asíncrona. Cada estado de consulta pertenece a un único usuario.
	\item Un \textbf{documento} puede contener cero o muchos \textbf{chunks}. La cardinalidad mínima es cero porque un documento puede encontrarse cargado en el sistema pero todavía no haber sido procesado ni fragmentado. Cada \eng{chunk} pertenece a un único documento.
	\item Cada \textbf{chunk} tiene asociado un único \textbf{embedding}, que almacena la información necesaria para gestionar su representación vectorial. Del mismo modo, cada \eng{embedding} pertenece a un único \eng{chunk}.
	\item La relación entre \textbf{consultas} y \textbf{chunks} es de muchos a muchos. Una consulta puede utilizar múltiples fragmentos documentales para generar una respuesta, mientras que un mismo fragmento puede participar en diferentes consultas realizadas por distintos usuarios. Esta relación se implementa mediante la entidad intermedia \textbf{ConsultaChunk}.
\end{itemize}

\begin{figure}
\centering
\begin{minipage}[c]{\textwidth}
    \centering
    \begin{subfigure}[t]{\textwidth}
        \centering
        \includegraphics[width=\textwidth]{Diagramas/Entidad_Relacion}
        \caption{Diagrama Entidad/Relación simplificado.}
    \end{subfigure}
\end{minipage}
\hfill
\begin{minipage}[c]{\textwidth}
    \centering
	\vspace*{0.2cm}
    \begin{subfigure}[t]{\textwidth}
        \centering
        \includegraphics[width=\textwidth]{Diagramas/herencia_tareas_asincronas}
        \caption{Diagrama Entidad/Relación de las tareas asíncronas.}
    \end{subfigure}
\end{minipage}
\caption{Diagramas Entidad/Relación del sistema.}
\label{fig:entidad_relacion}
\end{figure}

\section{Diseño arquitectónico} \label{dis:arc}
El objetivo del diseño arquitectónico es garantizar la modularidad, la escalabilidad, la mantenibilidad, la portabilidad y la reproducibilidad del sistema desarrollado. Es decir, se busca tener una separación clara de responsabilidades entre la interfaz, la lógica, la persistencia y los servicios de Inteligencia Artificial con una organización del código en componentes cohesivos (\eng{Blueprints}, servicios y modelos), pudiendo replicarse los servicios de manera sencilla. Además, sustituir los componentes debe tener un impacto mínimo en el resto del código. 

Este sistema se basa en una arquitectura cliente–servidor multicapa, combinada con una arquitectura \eng{RAG} (\eng{Retrieval-Augmented Generation}), en la capa de aplicación, estructurada en dos \eng{pipelines} diferenciados, el de indexación y el de generación. El despliegue se lleva a cabo utilizando contenedores \eng{Docker} para almacenar el código junto con sus dependencias.

\subsection{Arquitectura general del sistema y despliegue}
PythIA se ha diseñado como un sistema distribuido basado en servicios (\eng{service-oriented}~\cite{diseño:SOA}), donde cada componente desempeña una función específica dentro de la arquitectura global. Esta separación de responsabilidades facilita el mantenimiento, la escalabilidad y la sustitución de componentes sin afectar al resto del sistema. La arquitectura se compone de los siguientes elementos principales:
\begin{itemize}
	\item \textbf{Cliente \eng{web}}: se corresponde con el navegador utilizado por el usuario para acceder a la aplicaicón y realizar consultas.
	\item \textbf{Servidor Nginx}: actúa como punto de entrada al sistema (\eng{proxy} inverso), gestionando el protocolo de seguridad HTTPS a través del certificado SSL/TLS, y reenviando las peticiones hacia la aplicación.		
	\item \textbf{Capa de aplicación}: aplicación \eng{web} implementada mediante \eng{Flask} y desplegada sobre el servidor WSGI \eng{Gunicorn}. Constituye el núcleo del sistema y coordina la interacción entre los distintos servicios internos.
	\item \textbf{Base de datos relacional PostgreSQL}: almacena la información persistente de la aplicación, incluyendo usuarios, consultas, documentos y estados de las tareas asíncronas. 
	\item \textbf{Base de datos vectorial Qdrant}: almacena las representaciones vectoriales (\eng{embeddings}) utilizadas durante la recuperación semántica de información.	
	\item \textbf{Servicio de modelos \eng{LLM}}: implementado mediante Ollama, proporciona acceso a los modelos de lenguaje empleados para la generación de respuestas y la creación de \textit{embeddings}.
	\item \textbf{Contenedor de pruebas}: se utiliza para la ejecución automatizada de pruebas funcionales y de cobertura del sistema.
\end{itemize}

La figura~\ref{fig:Arquitectura/ComunicacionArquitectura.pdf}) muestra la organización general de los componentes y las comunicaciones existentes entre ellos. Todos los servicios se despliegan mediante contenedores Docker, lo que permite aislar dependencias, simplificar la instalación y garantizar la reproducibilidad del entorno de ejecución. El despliegue se define en el archivo \file{Docker Compose}. Este archivo coordina los seis contenedores que forman parte del sistema,  haciendo que cada componente puede evolucionar de manera independiente manteniendo una arquitectura modular y fácilmente portable entre distintos entornos de ejecución.

Los contenedores definidos en le archivo \file{Docker Compose} son:
\begin{enumerate}
	\item \file{db} $\to$ base de datos relacional.
	\item \file{qdrant} $\to$ base de datos vectorial.
	\item \file{ollama}, \file{ollama\_ocr} y \file{ollama-init} $\to$ servicio \eng{LLM}.
	\item \file{web} $\to$ aplicación \eng{web} Flask levantada usando Gunicorn.
	\item \file{nginx} $\to$ \eng{proxy} inverso y punto de terminación TLS~\cite{diseño:TLS}\footnote{TLS es un protocolo que asegura la privacidad y la seguridad de los datos} de la aplicación.
	\item \file{test} $\to$ ejecución de pruebas.
\end{enumerate}

Gracias a esta separación se garantiza la seguridad ya que no se expone directamente la aplicación \eng{Flask}, únicamente el servidor Nginx se encuentra expuesto al exterior. Por lo cual la comunicación entre componentes se realiza principalmente dentro de al red interna. Para llevar acabo esta comunicación los servicios se comunican entre sí mediante el nombre de servicio,es decir, el DNS interno, y se conectan a través de los puertos expuestos internamente.
 
El flujo general de una petición al sistema se muestra en la figura~\ref{fig:Arquitectura/Peticiones.pdf}. Cuando un usuario accede a la aplicación, la petición es recibida por Nginx y redirigida a la aplicación Flask ejecutada sobre Gunicorn. A partir de ese momento, la lógica de negocio interactúa con PostgreSQL, Qdrant y Ollama según las necesidades de cada operación, devolviendo finalmente la respuesta generada al usuario. 
 
\imagendiagram{Arquitectura/ComunicacionArquitectura.pdf}{Diagrama de la comunicación entre los principales componentes del sistema.}{1.2}

\imagendiagram{Arquitectura/Peticiones.pdf}{Diagrama de la gestión de peticiones del sistema.}{1.2}

\subsection{Arquitectura cliente-servidor multicapa}
Desde el punto de vista lógico el sistema se estructura siguiendo una arquitectura cliente-servidor multicapa~\cite{diseño:tresNiveles}, en la que las responsabilidades del sistema se distribuyen en diferentes niveles funcionales. Esta separación permite desacoplar la interfaz de usuario de la lógica de negocio y de los mecanismos de almacenamiento, facilitando el mantenimiento, la escalabilidad y la evolución del sistema.

La organización general de estas capas puede observarse en la figura~\ref{fig:Arquitectura/ArquitecturaLogica.pdf}. En ella se distinguen tres niveles claramente diferenciados. Estos niveles se corresponden con la capa de presentación, la capa de aplicación y la capa de persistencia. Cada una de ellas desempeña una función específica dentro del sistema. 

\imagendiagram{Arquitectura/ArquitecturaLogica.pdf}{Arquitectura cliente-servidor multicapa.}{1}

\begin{itemize}
	\item \textbf{Capa de presentación}: constituye la interfaz de usuario (IU). Está formada por las plantillas HTML generadas mediante Jinja2 y presentadas en el navegador. La aplicación sigue un enfoque SSR (\eng{Server-Side Rendering}), donde las páginas se renderizan en el servidor antes de enviarse al cliente. Este enfoque reduce la complejidad en el cliente, simplifica el control de acceso, mejora la trazabilidad y simplifica la integración de la lógica del sistema.
	\item \textbf{Capa de aplicación}: contiene la lógica funcional de la aplicación. Está formada por los controladores organizados mediante \textit{Blueprints}, los servicios de negocio y los procesos asociados al sistema \eng{RAG}. Esta capa recibe las peticiones procedentes de la interfaz, ejecuta las operaciones necesarias y coordina la comunicación con los sistemas de almacenamiento y los servicios de inteligencia artificial.
	\item \textbf{Capa de persistencia}: agrupa los mecanismos encargados del almacenamiento y recuperación de la información. Incluye tanto la base de datos relacional PostgreSQL, utilizada para la persistencia de usuarios, consultas, documentos y estados, como la base de datos vectorial Qdrant, empleada para almacenar los vectores utilizados durante la recuperación semántica. Asimismo, esta capa integra los servicios externos necesarios para el funcionamiento del sistema \eng{RAG}, como los modelos de lenguaje y generación de \eng{embeddings}.	
\end{itemize}

\subsection{Arquitectura de cuatro capas de la aplicación}

Aunque la arquitectura cliente-servidor permite comprender la organización general del sistema, resulta conveniente detallar con mayor precisión la estructura interna de la aplicación. Para ello, PythIA se ha diseñado siguiendo una arquitectura de cuatro capas (ver figura~\ref{fig:Arquitectura/Arquitectura_4capas.pdf})), en la que se separan claramente las responsabilidades asociadas a la interacción con el usuario, la lógica de negocio, la persistencia de la información y la integración con servicios externos.

\imagendiagram{Arquitectura/Arquitectura_4capas.pdf}{Arquitectura de 4 capas de la aplicación PythIA detallando la interacción entre los controladores, servicios y entidades.}{1}

\subsection{Arquitectura de la aplicación \eng{web}}
La aplicación \eng{web} sigue una arquitectura MVC (\eng{Model-View-Controller}). Esta arquitecturaes típica de las aplicaciones \eng{Flask} \eng{Server-Side-Rendered} (SSR)~\cite{diseño:SSR}. La técnica SSR consiste en generar el código de las páginas en el lado del servidor y entregar al cliente el código \file{HTML} completamente montado con toda la información que se debe mostrar. 

El modelo vista intención (MVC)~\cite{diseño:MVC} es un patrón arquitectónico de software que separa el código en tres capas independientes. La capa modelo donde se incluye la lógica y los datos , la capa vista donde se implementa la interfaz de usuario (\eng{frontend}) y la capa controlador que actúa de intermediario entre las acciones del usuario en las vistas y la lógica de la aplicación (rutas). En esta aplicación se implementa adaptado al ecosistema \eng{Flask} (ver diagrama~\ref{fig:Arquitectura/ModeloVistaControlador}).

\imagendiagram{Arquitectura/ModeloVistaControlador.pdf}{Diagrama de la arquitectura Modelo-Vista-Controlador (MVC) del sistema.}{1.2}

\begin{itemize}
	\item La \textbf{vista}: está compuesta por las platillas \file{HTML} renderizadas mediante el motor de plantillas \file{Jinja2}, responsables de presentar la interfaz y capturar las interacciones del usuario. Además, incorpora los recursos estáticos (\file{css} y \file{JavaScript}). Aunque, la aplicación incorpora \eng{JavaScript} para mejorar la experiencia de usuario, la generación principal del contenido es \eng{server-side}, por lo que la arquitectura sigue siendo fundamentalmente SSR.
	\item El \textbf{controlador}: se implementa mediante rutas \eng{Flask} organizadas en \eng{Blueprints}. El uso de \eng{Blueprints} permite mantener alta cohesión dentro de cada módulo y bajo acoplamiento entre funcionalidades. Se encarga de gestionar las peticiones HTTP, realizar validaciones básicas (como las de los formularios), así como controlar de acceso por rol. También, tiene la responsabilidad de invocar los servicios donde se encuentra la lógica, coordinar el \eng{pipeline} \eng{RAG} y generar las respuestas HTTP (renderizando la vista adecuada) . 
	\item El \textbf{modelo}: esta compuesto por las entidades de datos persistentes (Usuario, Consulta, Documento, \eng{Chunk}, \eng{Embedding} y la relación \eng{ConsultaChunk}) y por la lógica de negocio asociada (gestión de usuarios, administración documental, \eng{pipelines} de consulta e indexación \eng{RAG}, control de estados de procesos largos, etc). Se usa una base de datos relacional (\eng{PostgresSQL}) para asegurar la persistencia de los datos de la aplicación y una base de datos vectorial (\eng{Qdrant}) para la recuperación semántica. Esta capa permite desacoplar la lógica de negocio de la base de datos subyacente.
\end{itemize} 

Con estas capas se mantiene el desacoplamiento de la aplicación \eng{web}, haciendo que los controladores organicen el flujo para que la interfaz de usuario no implemente reglas de negocio y sea el modelo el que encapsule los datos y los procesos del domino. 

\subsection{Arquitectura del sistema \eng{RAG}}
El sistema \eng{RAG} se integra dentro de la capa de aplicación. Es el encargado de integrar en el sistema la generación de respuestas. u diseño se basa en la combinación de recuperación semántica sobre una base de datos vectorial y generación de texto mediante un modelo de lenguaje (\eng{LLM}). A nivel de arquitectura se divide en dos \eng{pipelines}~\cite{book:RAG_SimpleGuide}, el de indexación y el de generación (ver diagrama~\ref{fig:Arquitectura/ArquitecturaRAG}).

\imagendiagram{Arquitectura/ArquitecturaRAG.pdf}{Diagrama de la arquitectura del modelo \eng{RAG}.}{1.3}

El \textbf{\eng{pipeline} de indexación} se ejecuta cuando se cargan los documentos en la aplicación, ya bien sea de manera manual o mediante \eng{web scraping}. Por cada documento se va a llevar a cabo una conversión a \file{Markdown} para que todos tengan una estructura similar que el \eng{LLM} entienda fácilmente. De estos documentos se extrae, limpia y normaliza su texto, posteriormente se segmenta en \eng{chunks} con solapamiento y se generar sus \eng{embeddings}. Los cuales se van a almacenar en la base de datos vectorial (\eng{Qdrant}). Este proceso se ejecuta de forma asíncrona para evitar bloquear la interfaz \eng{web}. Cuando acaba la indexación de los documentos se envía un correo electrónico al usuario parta notificárselo.

El \textbf{\eng{pipeline} de generación} se ejecuta cuando el usuario hace una consulta. Su objetivo es dar una respuesta contextualizada utilizando los \eng{chunks} ya indexados en el sistema. Para ello va a generar los \eng{embeddings} de la consulta y a recuperar los \eng{chunks} más relevantes para generar la respuesta de \eng{Qdrant}. Los \eng{chunks} más relevantes se calculan por la similitud de sus \eng{embeddings} conlos de la consulta. Los usa para construir el \eng{prompt} aumentado, el cual envía al modelo \eng{LLM} para que responda a la consulta. Finalmente, muestra su respuesta al usuario y registra la consulta en la base de datos relacional. Este \eng{pipeline} combina recuperación semántica y generación de texto, asegurando trazabilidad al devolver los metadatos de los documentos fuente.

\section{Diseño procedimental} \label{dis:proc}

\section{Diseño de interfaces} \label{dis:int}

El diseño de las interfaces se ha desarrollado de manera incremental. Inicialmente se hicieron bocetos genéricos para plasmar las principales funcionalidades de la aplicación. Posteriormente se fue iterando sobre esos diseños hasta llegar a las interfaces finales de la \eng{web}.

Los prototipos iniciales que se crearon se ilustran en las figuras ~\ref{fig:Diseño_interfaces/prototipos/Inicio}, ~\ref{fig:Diseño_interfaces/prototipos/Login}, ~\ref{fig:Diseño_interfaces/prototipos/Singup}, ~\ref{fig:Diseño_interfaces/prototipos/Main}, ~\ref{fig:Diseño_interfaces/prototipos/Historial}, ~\ref{fig:Diseño_interfaces/prototipos/Main_details}, ~\ref{fig:Diseño_interfaces/prototipos/RAG}, ~\ref{fig:Diseño_interfaces/prototipos/Users}, ~\ref{fig:Diseño_interfaces/prototipos/Add_user}, ~\ref{fig:Diseño_interfaces/prototipos/Edit_user} y ~\ref{fig:Diseño_interfaces/prototipos/Documents}.

\imagenflotante{Diseño_interfaces/prototipos/Inicio}{Prototipo de la página de inicio.}
\imagenflotante{Diseño_interfaces/prototipos/Login}{Prototipo de la página de inicio de sesión.}
\imagenflotante{Diseño_interfaces/prototipos/Singup}{Prototipo de la página de registro.}
\imagenflotante{Diseño_interfaces/prototipos/Main}{Prototipo de la página principal.}
\imagenflotante{Diseño_interfaces/prototipos/Historial}{Prototipo de la página del historial.}
\imagenflotante{Diseño_interfaces/prototipos/Main_details}{Prototipo de los detalles del historial.}
\imagenflotante{Diseño_interfaces/prototipos/RAG}{Prototipo de la página de consulta.}
\imagenflotante{Diseño_interfaces/prototipos/Users}{Prototipo de la página de gestión de usuarios.}
\imagenflotante{Diseño_interfaces/prototipos/Add_user}{Prototipo de la página de añadir usuarios.}
\imagenflotante{Diseño_interfaces/prototipos/Edit_user}{Prototipo de la página de edición del usuario.}
\imagenflotante{Diseño_interfaces/prototipos/Documents}{Prototipo de la página de gestión de documentos.}

Aunque estos prototipos se asemejan mucho a las interfaces finales, se han realizado diversos cambios. Por ejemplo, se ha cambiado la cabecera para aportar más información como el nombre del usuario que inicia sesión, su rol (usuario normal o administrador) y su último inicio de sesión (ver imagen de la cabecera final~\ref{fig:Diseño_interfaces/Cabecera}). Además, se ha añadido a la tabla del historial de consultas y a la de documentos cargados información relevante que no se había valorado inicialmente. En la tabla del historial de consultas se han añadido diversos campos para aportar toda la información relevante sobre la consulta. Añadiendo el documento al que pertenece el \eng{chunk} principal y la fecha en la que el usuario hizo la consulta. Para el administrador, también, se ha añadido una columna en la que se indica el nombre del usuario que hizo la consulta. Además, para mejorar la legibilidad de la tabla se ha decidido ubicar la respuesta debajo de la fila en un desplegable (ver imagen del historial de consulta definitivo~\ref{fig:Diseño_interfaces/Table_asks}). Por su parte, en la tabla donde se muestran los documentos cargados del sistema se han añadido dos columnas, una que indica el tamaño del documento y otra que indica el estado de ese documento en el sistema (cargado, indexado o fallido) (ver imagen del diseño final de la tabla documentos~\ref{fig:Diseño_interfaces/Table_documents}). Otro aspecto relevante es que en la página para realizar la consulta se ha añadido un desplegable para ver los \textit{chunks} utilizados (ver imagen de la página de consultas definitiva~\ref{fig:Diseño_interfaces/RAG_def}).

\imagenflotante{Diseño_interfaces/Cabecera}{Diseño de interfaz -- Cabecera.}

\imagenflotante{Diseño_interfaces/Table_asks}{Diseño de interfaz -- Historial de consultas.}

\imagenflotante{Diseño_interfaces/Table_documents}{Diseño de interfaz -- Tabla de documentos cargados.}

\imagenflotante{Diseño_interfaces/RAG_def}{Diseño de interfaz -- Interfaz de consulta al modelo.}

En estos diagramas  ~\ref{fig:Diagramas/Navegabilidad_admin} y ~\ref{fig:Diagramas/Navegabilidad_normal} se muestra como sería la navegabilidad de los usuario por la aplicación dependiendo del rol que tienen (administrador o usuario normal).
\begin{landscape}
	\imagendiagram{Diagramas/Navegabilidad_admin}{Diagrama de navegabilidad por la aplicación para el administrador.}{1.2}
	\imagendiagram{Diagramas/Navegabilidad_normal}{Diagrama de navegabilidad por la aplicación para el usuario normal.}{1.2}
\end{landscape}




